\subsection{Risks Management}
Possible risks affecting project progress are outlined in Section 3.3. 
With the timeline now established, contingency plans are defined to 
mitigate potential delays. This section presents each identified risk 
along with its probability, impact, and corresponding mitigation strategy.

\subsubsection{Coding Bugs}

\hfill
\textbf {Probability:} High

\hfill
\textbf{Impact:} Medium


Coding bugs are inevitable, making thorough testing essential. As 
outlined in Section 5.4, every code change undergoes rigorous testing 
across all use cases, including edge cases. In addition, during the 
Dev1 task, Visual Studio Code (VS Code) extensions such as ESLint and 
SonarQube are installed to enable automated bug detection from the very 
beginning of implementation. These extensions highlight coding issues 
directly in the editor, preventing the developer from running the 
application until the errors are resolved.

If bugs persist, additional debugging time is allocated. However, thanks 
to the installed extensions and consistent testing, the number of bugs 
is minimized, reducing the likelihood of significant delays. Furthermore, 
complex tasks already include built-in buffer time, and the reserve hours 
are available to cover any remaining delays that may arise.

\subsubsection{Tight Project Deadlines}

\hfill
\textbf{Probability:} High

\hfill
\textbf{Impact:} Low

Software development projects often struggle with tight deadlines. To 
overcome this challenge, the remaining 38 hours of overhead serve as a 
lifesaver. This can be particularly useful throughout the project, as 
unexpected delays are quite common.

\subsubsection{Lack of Experience in Certain Technologies}

\hfill
\textbf{Probability:} High

\hfill
\textbf{Impact:} Medium

Due to the author's limited coding experience, several technologies used 
in the application implementation fall outside their core expertise. Task 
Dev0 addresses this challenge by allocating time for self-study, helping 
to minimize potential delays. If significant delays still persist, 
assistance from the FSP team serves as a fallback. Seeking timely 
assistance is an essential software development skill that helps avoid 
unnecessary deadline extensions.

\subsubsection{Company's licenses and security permissions}

\hfill
\textbf{Probability:} High

\hfill
\textbf{Impact:} Low

Another risk involves FSP's security system. Cyberattacks are taken more 
seriously every day, often requiring developers to request and wait for 
permission to install resources and extensions. Consequently, new 
restrictions may limit initial resource availability. As a contingency, 
tasks with no pending dependencies, such as documentation, can be 
prioritized in such scenarios.
